\chapter{Addition: First Steps}\label{ch:g1:addition}

Three red marbles and two blue marbles: together, how many? Putting
together and counting the whole is \emph{adding} --- the first
operation, with its two famous signs $+$ and $=$.

\section{Putting together}

\begin{definition}[Addition]\label{def:g1:addition:def}
\emph{Adding} means putting together and counting everything. We
write
\[
3 + 2 = 5
\]
and read: ``three plus two equals five''. The answer of an addition
is called the \emph{sum}.
\end{definition}

\begin{omfigure}
\begin{tikzpicture}[scale=0.6]
  \foreach \i in {0,1,2} \fill[omThm] (\i*1.0,0) circle (0.3);
  \node[font=\large] at (3.0,0) {$+$};
  \foreach \i in {0,1} \fill[omDef] (3.8+\i*1.0,0) circle (0.3);
  \node[font=\large] at (5.8,0) {$=$};
  \foreach \i in {0,1,2} \fill[omThm] (6.6+\i*1.0,0) circle (0.3);
  \foreach \i in {0,1} \fill[omDef] (9.6+\i*1.0,0) circle (0.3);
  \node[below, font=\small] at (5.6,-0.7)
    {$3 + 2 = 5$: three and two make five};
\end{tikzpicture}

{\small An addition in pictures: the two groups together make one
group of five.}
\end{omfigure}

\begin{example}[The order does not matter]\label{ex:g1:addition:order}
$2 + 3$ and $3 + 2$ put together the same marbles: both make $5$.
To compute $2 + 9$, it is faster to start from the bigger number:
say ``nine'', then count on two --- ``ten, eleven''. So
$2 + 9 = 11$.
\end{example}

\section{The ten friends}

\begin{example}[Complements to 10]\label{ex:g1:addition:friends}
Two numbers that make $10$ together are ``ten friends'':
\[
1 + 9, \quad 2 + 8, \quad 3 + 7, \quad 4 + 6, \quad 5 + 5 .
\]
Ten fingers show them all: raise $3$ fingers, and the $7$ folded
ones are the friend. Knowing the ten friends by heart makes many
additions fast.
\end{example}

\begin{omfigure}
\begin{tikzpicture}[scale=0.72]
  % ten-frame: 2 rows of 5 cells
  \draw[thick] (0,0) grid[step=1] (5,2);
  \draw[very thick] (0,0) rectangle (5,2);
  \foreach \i in {0,1,2} \fill[omDef] (\i+0.5,1.5) circle (0.3);
  \foreach \i in {3,4} \draw[omThm, thick] (\i+0.5,1.5) circle (0.3);
  \foreach \i in {0,...,4} \draw[omThm, thick] (\i+0.5,0.5) circle (0.3);
  \node[right, font=\small, align=left] at (5.6,1.0)
    {a \emph{ten-frame}: $10$ cells.\\
     $3$ filled, $7$ empty: $3 + 7 = 10$};
\end{tikzpicture}

{\small The ten-frame shows every pair of ten friends at a glance:
what is filled and what is missing always make $10$ together.}
\end{omfigure}

\begin{method}[Crossing the ten]\label{met:g1:addition:crossten}
To compute $8 + 5$:
\begin{enumerate}
  \item take from the $5$ what the $8$ needs to reach $10$: that is
        $2$ (its ten friend);
  \item $8 + 2 = 10$, and $3$ remains from the $5$;
  \item $10 + 3 = 13$. So $8 + 5 = 13$.
\end{enumerate}
Passing through $10$ turns a hard sum into two easy ones.
\end{method}

\begin{omfigure}
\begin{tikzpicture}[scale=0.62]
  % ten-frame with 8 blue, 2 red (borrowed from the 5)
  \draw[thick] (0,0) grid[step=1] (5,2);
  \draw[very thick] (0,0) rectangle (5,2);
  \foreach \i in {0,...,4} \fill[omDef] (\i+0.5,1.5) circle (0.3);
  \foreach \i in {0,1,2} \fill[omDef] (\i+0.5,0.5) circle (0.3);
  \foreach \i in {3,4} \fill[omThm] (\i+0.5,0.5) circle (0.3);
  % 3 red left outside the frame
  \foreach \i in {0,1,2} \fill[omThm] (6.2+\i,1.0) circle (0.3);
  \node[below, font=\small] at (2.5,-0.25) {the frame is full: $10$};
  \node[below, font=\small] at (7.2,0.55) {$3$ spill over};
\end{tikzpicture}

{\small $8 + 5$ on a ten-frame: the $8$ blue counters take $2$ of the
$5$ red ones to fill the frame --- one full ten and $3$ more, so
$13$.}
\end{omfigure}

\begin{omfigure}
\begin{tikzpicture}[scale=0.6]
  \draw[-Stealth] (-0.4,0) -- (15.4,0);
  \foreach \i in {0,...,15} \draw (\i,-0.1) -- (\i,0.1);
  \foreach \i in {0,5,8,10,13,15}
    \node[below, font=\footnotesize] at (\i,-0.12) {\i};
  \draw[-Stealth, omDef, thick] (8,0.3) .. controls (9,0.95) ..
    (10,0.3);
  \node[omDef, font=\small] at (9,1.05) {$+2$};
  \draw[-Stealth, omThm, thick] (10,0.3) .. controls (11.5,1.15) ..
    (13,0.3);
  \node[omThm, font=\small] at (11.5,1.2) {$+3$};
\end{tikzpicture}

{\small $8 + 5$ on the number line: hop to $10$ first, then the rest.}
\end{omfigure}

\begin{example}[Doubles]\label{ex:g1:addition:doubles}
The doubles are worth learning by heart:
\[
1+1=2, \quad 2+2=4, \quad 3+3=6, \quad 4+4=8, \quad 5+5=10,
\]
\[
6+6=12, \quad 7+7=14, \quad 8+8=16, \quad 9+9=18, \quad 10+10=20 .
\]
Near-doubles come free: $6 + 7$ is one more than $6 + 6$, so $13$.
\end{example}

\section{Little problems}

\begin{example}\label{ex:g1:addition:problem}
``Lea has $7$ stickers. She gets $6$ more. How many now?''
Getting more means adding: $7 + 6 = 13$ (through ten:
$7 + 3 = 10$, then $+3$). Lea has $13$ stickers. Always finish with
a little sentence!
\end{example}

\section{Exercises}

\begin{exercise}[$\star$]\label{exo:g1:addition:1}
Draw the marbles and compute: $4 + 3$; \quad $5 + 2$; \quad
$6 + 4$.
\end{exercise}

\begin{exercise}[$\star$]\label{exo:g1:addition:2}
Compute by starting from the bigger number: $3 + 8$; \quad
$2 + 12$; \quad $4 + 9$.
\end{exercise}

\begin{exercise}[$\star$]\label{exo:g1:addition:3}
Give the ten friend of: $7$; \ $4$; \ $9$; \ $5$; \ $1$.
\end{exercise}

\begin{exercise}[$\star$]\label{exo:g1:addition:4}
Compute through the ten (\cref{met:g1:addition:crossten}):
$9 + 4$; \quad $8 + 7$; \quad $6 + 5$.
\end{exercise}

\begin{exercise}[$\star$]\label{exo:g1:addition:5}
Recite the doubles, then compute using them: $7 + 8$; \quad
$5 + 6$; \quad $9 + 8$.
\end{exercise}

\begin{exercise}[$\star$]\label{exo:g1:addition:6}
Complete the holes:
\[
6 + \;?\; = 10, \qquad
\;?\; + 5 = 12, \qquad
10 + \;?\; = 17 .
\]
\end{exercise}

\begin{exercise}[$\star$]\label{exo:g1:addition:7}
``In the bus there are $8$ children; $5$ grown-ups climb in. How
many people are in the bus?'' Write the addition and the answer
sentence.
\end{exercise}

\begin{exercise}[$\star$]\label{exo:g1:addition:8}
Tom has $6$ red cars and $7$ blue cars. Mia has $10$ cars. Who has
more cars?
\end{exercise}

\begin{exercise}[$\star\star$]\label{exo:g1:addition:9}
Three dice show $4$, $6$ and $5$. What is the total? (Add two dice
first --- choose the two that make a ten friend!)
\end{exercise}
